\section{The signed native certificate lattice}
\label{sec:tpc68-certificate-lattice}

The preceding theorems take absolute values at different stages.  Their
conclusions can be compared only after each is evaluated on the same
complete literal matrix.

\begin{definition}[Available certificates]
\label{def:tpc68-certificate-family}
For the gauged alias operator \(H\), define
\begin{align}
 C_{\rm star}
 &:=\sqrt{\Delta_H}\,s_*(H),
 \label{eq:tpc68-C-star}\\
 C_{\rm pair}
 &:=\rho(P),
 \label{eq:tpc68-C-pair}\\
 C_{{\rm walk},q}
 &:=\Omega_{2q}(H)^{1/(2q)},
 \label{eq:tpc68-C-walk}\\
 C_{{\rm mom},q}
 &:=\bigl(\tr H^{2q}\bigr)^{1/(2q)}
 \qquad(q\ge1).
 \label{eq:tpc68-C-moment}
\end{align}
Here \(P\) may be the optimal literal pairing-defect matrix or any other
proved symmetric nonnegative pair majorant.  Put
\begin{equation}
 \mathcal C_X
 :=
 \min\left\{
 C_{\rm star},C_{\rm pair},
 \inf_{q\ge1}C_{{\rm walk},q},
 \inf_{q\ge1}C_{{\rm mom},q}
 \right\}.
 \label{eq:tpc68-master-certificate}
\end{equation}
\end{definition}

At a fixed finite \(X\), all terms are finite and computable in principle.
The infima over \(q\) are mathematical certificates; their exact
evaluation by arbitrarily high powers is not claimed to be computationally
efficient.

\begin{theorem}[Exact native eigenvalue certificate]
\label{thm:tpc68-native-certificate-lattice}
For the atomic-normalized literal missing-native Gram \(K=I+H\),
\begin{align}
 \kappa_M
 &\le1+\mathcal C_X,
 \label{eq:tpc68-kappa-lattice}\\
 \alpha_M
 &\ge\max\{0,1-\mathcal C_X\}.
 \label{eq:tpc68-alpha-lattice}
\end{align}
Every individual term in
\eqref{eq:tpc68-master-certificate} may replace \(\mathcal C_X\).
\end{theorem}

\begin{proof}
\Cref{thm:tpc68-star-sandwich,thm:tpc68-Perron-certificate,thm:tpc68-even-walk-certificate,thm:tpc68-even-moment}
shows that \(\|H\|\) is at most every listed certificate.  Hence
\(\|H\|\le\mathcal C_X\).  Since
\[
 \lambda_{\max}(H)\le\|H\|,
 \qquad
 \lambda_{\min}(H)\ge-\|H\|,
\]
the eigenvalue identities
\eqref{eq:tpc68-upper-native-eigenvalue}--\eqref{eq:tpc68-lower-native-eigenvalue}
give the two bounds.  Positivity of \(K\) supplies the additional zero in
\eqref{eq:tpc68-alpha-lattice}.
\end{proof}

\begin{remark}[Exact but not automatically small]
\label{rem:tpc68-certificate-not-saving}
The theorem is an exact L0 implication.  It does not say that
\(\mathcal C_X=o(1)\), \(X^{o(1)}\), or even bounded.  Computing a large
certificate exactly is diagnostically useful but does not prove
cancellation.
\end{remark}

\subsection{Full-support components}

Let the actual collision graph of \(H\) have connected components
\(\mathscr C_X\).  After a simultaneous permutation of rows,
\[
 H=\bigoplus_{C\in\mathscr C_X}H_C,
 \qquad
 K=\bigoplus_{C\in\mathscr C_X}(I_C+H_C).
\]

\begin{proposition}[Worst-component rule]
\label{prop:tpc68-worst-component}
One has
\begin{align}
 \|H\|&=\max_{C\in\mathscr C_X}\|H_C\|,
 \label{eq:tpc68-component-norm}\\
 \kappa_M&=\max_{C\in\mathscr C_X}
 \lambda_{\max}(I_C+H_C),
 \label{eq:tpc68-component-upper}\\
 \alpha_M&=\min_{C\in\mathscr C_X}
 \lambda_{\min}(I_C+H_C).
 \label{eq:tpc68-component-lower}
\end{align}
Thus a uniform theorem must include every actual component.  Selecting a
favorable component or deleting a difficult row changes the operator.
\end{proposition}

\begin{proof}
These are the standard norm and extremal-eigenvalue formulas for a finite
orthogonal direct sum.
\end{proof}

\subsection{Quantifier-safe finite certificates}

The pairing \(\tau\), Abel ordering, Perron weight \(a\), power \(q\), and
component decomposition may all be chosen after the finite literal matrix
is known.  They still yield coefficient-uniform operator inequalities
because none depends on the vector to which \(H\) is later applied.  What
is forbidden is to choose or alter the physical carrier, erase extra
output labels, delete active rows, or rebuild the certificate separately
for each later coefficient vector while advertising one uniform operator
bound.
